\clearpage\begingroup\NoSourceHead\SourceLeaf{545b}
\renewcommand{\thefootnote}{\arabic{footnote})}\setcounter{footnote}{4}
\ArithmeticHead{88}{L}{BOLZA, GAUSS UND DIE VARIATIONSRECHNUNG.}
\begin{center}\textls[100]{3. Über den homogenen Körper grösster Attraktion}\\
\textls[100]{bei gegebener Masse.}\\
\small (Handbuch 16, Bb, S. 107, Werke, Bd. XI 1.)\end{center}
\noindent\hspace*{\parindent}\textsc{Gauss} gibt hier in gedrängter Kürze mit nur wenig Beweisandeutungen
die Lösungen der folgenden \textls[120]{vier Aufgaben}:
\begin{list}{\arabic{enumi})}{\usecounter{enumi}\setlength{\leftmargin}{2.4em}\setlength{\itemsep}{0pt}\setlength{\topsep}{0pt}}
\item Bestimmung des homogenen Körpers von gegebener Masse und grösster
Anziehung auf einen gegebenen Punkt in gegebener Richtung unter
Voraussetzung des \textsc{Newton}schen Anziehungsgesetzes;
\item Vergleichung der Anziehung dieses Körpers mit derjenigen einer Kugel
von gleicher Masse;
\item Bestimmung des Sphäroids von gegebener Masse und grösster Anziehung auf einen seiner Pole;
\item Vergleichung der Anziehung dieses Sphäroids mit derjenigen einer
Kugel von der gleichen Masse.
\end{list}
\par\noindent\hspace*{\parindent}Am Schluss bemerkt \textsc{Gauss}: »Eine ähnliche Untersuchung soll von Herrn
\textsc{Playfair} im VI. Band der Transactions of the Royal Society of Edinburgh
angestellt sein (1812)«. Diese Bemerkung stützt sich auf einen ganz kurzen
Bericht über die Arbeit von J. \textsc{Playfair} in den Göttinger Gelehrten Anzeigen
\clearpage\SourceLeaf{546}
\ArithmeticHead{89}{R}{VEREINZELTE ANWENDUNGEN DER VARIATIONSRECHNUNG.}
\noindent für 1818, S. 860, aus dem wenig mehr zu ersehen war, als was der Titel der
\textsc{Playfair}schen Arbeit, (die übrigens schon im Jahr 1807 vor der Gesellschaft
gelesen worden war), besagte: \textit{Of the Solids of Greatest Attraction, or those
which, among all the Solids that have certain Properties, attract with the greatest
Force in a given Direction.} \textsc{Gauss}’ Vermutung war richtig: \textsc{Playfair} beschäftigt
sich in der Tat in der angegebenen Arbeit unter anderem mit den drei ersten
der oben aufgeführten Aufgaben und löst sie, wie wir sehen werden, — mit
einer Ausnahme — richtig. —
\par\noindent\hspace*{\parindent}Die \textls[120]{erste Aufgabe} erledigt \textsc{Gauss} durch die lakonische Bemerkung:
»Man überzeugt sich leicht, dass jedes Teilchen an der Oberfläche des Körpers
grösster Anziehung gleich starke Anziehung beitragen muss. Daher entsteht
der Körper durch Umdrehung einer Kurve, deren Gleichung ist
\begin{equation*}\tag{6}
\frac{\cos u}{rr}=A\text{\textsuperscript{5)}«}.
\end{equation*}
\setcounter{footnote}{5}\footnotetext{Um Übereinstimmung mit den bei \textsc{Gauss} weiter folgenden Formeln für die Masse und die Anziehung dieses Körpers herzustellen, müsste hier $\frac1{A^2}$ statt $A$ geschrieben werden.}
Dabei sind $r,u$ Polarkoordinaten mit dem gegebenen Punkt als Pol und der
gegebenen Richtung als Achse.
\par\noindent\hspace*{\parindent}Gemeint ist hier ohne Zweifel dieselbe Überlegung, durch die bereits
R. J. \textsc{Boscovich}\footnote{\textit{Problema mecanicum de solido maximae attractionis}, Memorie sopra la fisica e istoria naturale di
diversi valentuomini, T. I, Lucca 1743, S. 63—88 (auf der Landesbibliothek in Stuttgart vorhanden). Die
Aufgabe war \textsc{Boscovich} zwei Jahre vorher von \textsc{De Montigni} gestellt worden, und zwar in folgender
Fassung: »Data quantitate materiae punctum attrahentis, in quacunque lege distantiarum invenire solidum
ipsam continens, quod maxime omnium attrahat ipsum punctum positum in axe solidi producto ad datam
distantiam ab ipsius solidi vertice propiore«. \textsc{Boscovich} gibt an, dass er die im Text geschilderte »geometrische« Lösung ohne Mühe gefunden habe, dass ihm dagegen die Lösung derselben Aufgabe mittels
Integralrechnung grosse Schwierigkeiten bereitet habe. Auch diese zweite »analytische« Lösung, die auf
\textsc{Jakob Bernoulli}s Methode der Variation zweier benachbarter Ordinaten der Meridiankurve gegründet ist,
zeichnet sich durch grosse Einfachheit und Eleganz aus und zeigt deutlich den inneren Zusammenhang
zwischen der geometrischen und der analytischen Lösung. Am Schluss seiner Arbeit hebt \textsc{Boscovich} als
Vorzug der geometrischen Methode vor der analytischen hervor, dass die erstere auch auf den Fall anwendbar sei, wo die zur Vergleichung herangezogenen Körper nicht auf Rotationskörper beschränkt werden.} (1743) die Aufgabe zunächst unter Beschränkung auf Rotationskörper in etwas modifizierter Formulierung und unter Annahme einer
beliebigen, nur von der Entfernung abhängigen Anziehungskraft gelöst hatte:
Angenommen die Anziehungskomponente in der Richtung der Achse wäre in
\par\noindent{\footnotesize X 2 Abh. 5.\hfill 12}\par
\clearpage\SourceLeaf{547}
\ArithmeticHead{90}{L}{BOLZA, GAUSS UND DIE VARIATIONSRECHNUNG.}
\noindent zwei Punkten der Oberfläche des Körpers grösster Anziehung verschieden, so
schneide man aus dem Körper ein kleines Stück in der Umgebung des Punktes
geringerer Anziehung heraus und lege es in der Umgebung des Punktes grösserer
Anziehung von aussen an den Körper an, so würde man einen Körper gleicher
Masse, aber grösserer Anziehung erhalten. \textsc{Boscovich} knüpft hieran zugleich
einen Hinlänglichkeitsbeweis, indem er bemerkt, dass im Falle, wo die Anziehungskraft mit abnehmender Entfernung zunimmt, die Anziehung auf den
gegebenen Punkt in der Richtung der Achse für Punkte im Innern des gefundenen Körpers grösser, für diejenigen ausserhalb kleiner ist als für Punkte
auf der Oberfläche, woraus folgt, dass jede Veränderung des Körpers bei gleich
bleibender Masse die Anziehung vermindert.
\par\noindent\hspace*{\parindent}Für den speziellen Fall des \textsc{Newton}schen Anziehungsgesetzes findet sich
der erste Teil dieses Schlusses wieder bei \textsc{Saint Jacques de Silvabelle}\footnote{\textit{Problème: Supposant la loi d’attraction en raison inverse du carré de la distance, trouver la nature
du solide de la plus grande attraction}, Mémoires de Mathématiques et de Physiques, presentés à l’Académie
Royale des Sciences par divers Sçavans, Tome I, Paris 1750, S. XVII und 175/176. Die Arbeit ist datiert
7. Juli 1745.} (1745),
die beiden Teile bei J. \textsc{Playfair} (1807), der zweite Teil bei K. H. \textsc{Schellbach}\footnote{\textit{Probleme der Variationsrechnung}, Journal für die reine und angewandte Mathematik, Bd. 41 (1851),
S. 343—345.} (1851). Alle diese Autoren, \textsc{Gauss} inbegriffen, scheinen unabhängig
von einander auf diese Lösung des Problems ohne Variationskalkül gekommen
zu sein.
\par\noindent\hspace*{\parindent}Schliesslich hat neuerdings \textsc{Spijker}\footnote{\textit{Der Körper grösster Anziehung eines Ellipsoides}, Dissertation, Zürich 1904.} denselben Gedanken verwertet, um
allgemeiner ohne Benutzung der Variationsrechnung zu beweisen, dass die
Begrenzung desjenigen Körpers, der, als Integrationsbereich genommen, das
Integral $\iiint F_1(x,y,z)\,dx\,dy\,dz$ zu einem Extremum macht, während gleichzeitig das Integral $\iiint F_2(x,y,z)\,dx\,dy\,dz$ einen gegebenen Wert haben soll,
der Gleichung genügen muss
\[
F_1(x,y,z)-kF_2(x,y,z)=0,
\]
wo $k$ eine Konstante ist.
\par\noindent\hspace*{\parindent}Nachdem die Meridiankurve des Körpers grösster Anziehung gefunden
ist, bietet die \textls[120]{zweite Aufgabe}, die Berechnung des Verhältnisses der An-
\clearpage\SourceLeaf{548}
\ArithmeticHead{91}{R}{VEREINZELTE ANWENDUNGEN DER VARIATIONSRECHNUNG.}
\noindent ziehung dieses Körpers auf den gegebenen Punkt zur Anziehung einer Kugel
von gleicher Masse auf einen Punkt ihrer Oberfläche, keine Schwierigkeiten.
\textsc{Gauss} findet für das Verhältnis den Wert $3:\sqrt[3]{25}=1{,}025985$. Dasselbe Resultat gibt er auch in einer Fussnote zur Einleitung seiner \textit{Principia generalia
theoriae figurae fluidorum in statu aequilibrii}\footnote{Werke, Bd. V, S. 31.} (1829), und zwar ausdrücklich als
ein bekanntes Resultat, aber ohne Literaturnachweisung. Das bezieht sich
ohne Zweifel darauf, dass \textsc{Playfair}\footnote{loc. cit., S. 198.} in der oben zitierten Arbeit, die wohl
inzwischen \textsc{Gauss} zu Gesicht gekommen war, dasselbe Resultat abgeleitet
hatte. Später ist derselbe Satz nochmals von \textsc{Schellbach}\footnote{loc. cit. 8), S. 345.} bewiesen worden
und ist dann in die Lehrbücher übergegangen\footnote{\textsc{Moigno-Lindelöf}, \textit{Calcul des variations}, Paris 1861, S. 247.}. —
\par\noindent\hspace*{\parindent}Um die \textls[120]{dritte Aufgabe} zu lösen, die Bestimmung des Sphäroids von
gegebener Masse und grösster Anziehung im Pol, bezeichnet \textsc{Gauss} das Verhältnis der kleinen Achse des Sphäroids zur grossen Achse mit $\cos\varphi$ und berechnet dann das Verhältnis der Anziehung des Sphäroids im Pol zur Anziehung einer Kugel von gleicher Masse auf einen Punkt ihrer Oberfläche;
er findet dafür den Wert
\begin{equation*}\tag{7}
\frac{3\cos\varphi^{\frac53}}{\sin\varphi^3}(\operatorname{tg}\varphi-\varphi).
\end{equation*}
»Dieser Ausdruck wird ein Maximum für
\begin{equation*}\tag{8}
\varphi=\frac{9\operatorname{tg}\varphi+2\operatorname{tg}\varphi^3}{9+5\operatorname{tg}\varphi^2},
\end{equation*}
also für $\varphi=43^0\,59'\,2''$.« Hieraus berechnet dann \textsc{Gauss} noch das Verhältnis
der beiden Anziehungen zu $1{,}02204$ (\textls[120]{vierte Aufgabe}).
\par\noindent\hspace*{\parindent}Mit derselben Aufgabe hatte sich auch bereits \textsc{Playfair}\footnote{loc. cit., S. 220—225.} beschäftigt: Er
führt ebenfalls den Winkel $\varphi$ ein, berechnet dann die Anziehung des Sphäroids von der gegebenen Masse $\frac{4\pi n^3}3$ auf einen der Pole, die sich von dem
\textsc{Gauss}schen Ausdruck (7), wie a priori klar ist, nur um den konstanten Faktor
$\frac{4\pi n}3$ unterscheidet, und wird so zu derselben Gleichung (8) für $\varphi$ geführt wie
\par\noindent\hfill{\footnotesize $12^{*}$}\par
\clearpage\SourceLeaf{549}
\ArithmeticHead{92}{L}{BOLZA, GAUSS UND DIE VARIATIONSRECHNUNG.}
\noindent \textsc{Gauss}. Hier begnügt er sich nun aber, ohne zu untersuchen, ob die Gleichung noch weitere Wurzeln besitzt, mit der auf der Hand liegenden Lösung
$\varphi=0$, von der er, ohne auf eine Untersuchung der zweiten Ableitung einzugehen, behauptet, dass sie das Maximum liefert. So gelangt er im Gegensatz zu \textsc{Gauss} zu dem Resultat, dass die Kugel unter allen Sphäroiden von
gleicher Masse die grösste Anziehung auf einen Punkt im Pol ausübt.
\par\noindent\hspace*{\parindent}Mit Rücksicht auf diesen Widerspruch mag hier eine Diskussion der
Maximalaufgabe Platz finden. Führt man mit \textsc{Playfair} statt des Winkels $\varphi$
die Grösse $t=\operatorname{tg}\varphi$ ein, so geht der Ausdruck für die Anziehung des Sphäroids von der Masse $\frac{4\pi n^3}3$ in einem seiner Pole über in
\[
F(t)=\frac{4\pi n(1+t^2)^{\frac23}(t-\operatorname{Arc\,tg}t)}{t^3}.
\]
Für die Ableitung ergibt sich hieraus
\[
F'(t)=\frac{4\pi n}3\,\frac{(9+5t^2)\operatorname{Arc\,tg}t-(9t+2t^3)}{t^4(1+t^2)^{\frac13}}
\]
und daraus als Bedingung für ein Extremum
\[
f(t)\equiv\operatorname{Arc\,tg}t-\frac{9t+2t^3}{9+5t^2}=0
\]
in Übereinstimmung mit der Gleichung (8) von \textsc{Gauss} und \textsc{Playfair}. Da nun
\[
f'(t)=\frac{2t^4(3-5t^2)}{(1+t^2)(9+5t^2)^2},
\]
so schliesst man, dass während $t$ von $0$ bis $\infty$ wächst, die Funktion $f(t)$ von
$0$ anfangend bis zu einem positiven, für $t=\sqrt{\frac35}$ erreichten Maximalwert zunimmt und von da an bis $-\infty$ beständig abnimmt, also noch einmal verschwinden muss für einen Wert $t_0$, der etwas kleiner als $1$ sein muss, da
$f(1)=-0{,}0003$. Hieraus folgt dann, dass die Funktion $F(t)$ für $t=0$ ein
(relatives) Minimum, für $t=t_0$ das (absolute) Maximum besitzt, in Übereinstimmung mit \textsc{Gauss}, während sich das Resultat von \textsc{Playfair} als unrichtig
erweist.
\par\noindent\hspace*{\parindent}Ganz analoge Formeln gelten für das verlängerte Rotationsellipsoid. Setzt
man hier das Verhältnis der Polarachse zur Äquatorialachse gleich $\cosh\varphi$, so
erhält man für die Anziehung des verlängerten Rotationsellipsoids von der
\clearpage\SourceLeaf{550}
\ArithmeticHead{93}{R}{VEREINZELTE ANWENDUNGEN DER VARIATIONSRECHNUNG.}
\noindent Masse $\frac{4\pi n^3}3$ im Pol den Ausdruck
\[
F=\frac{4\pi n\cosh\varphi^{\frac53}(\varphi-\tanh\varphi)}{\sinh\varphi^3},
\]
aus dem bei ganz analoger Behandlung folgt, dass die Anziehung beständig
abnimmt, während $\varphi$ von $0$ bis $\infty$ wächst, ein Resultat, das \textsc{Playfair} am
Schluss seiner Untersuchung als Vermutung ausspricht.
\begin{center}\small SCHLUSSBEMERKUNG.\end{center}
{\small\noindent\hspace*{\parindent}Bei der Abfassung des vorstehenden Aufsatzes habe ich mich in weitgehendem Masse der Unterstützung von F. \textsc{Klein}, A. \textsc{Kneser}, G. \textsc{Prange}, L. \textsc{Schlesinger} und P. \textsc{Stäckel} ($\dagger$) zu erfreuen gehabt,
die teils durch Mitlesen der Korrekturen, teils durch wertvolle Literaturnachweise und kritische Bemerkungen
die Arbeit wesentlich gefördert haben. Allen diesen Herren sei auch an dieser Stelle mein herzlichster Dank
ausgesprochen.\hfill\textsc{Bolza}.\par}
\begin{center}\rule{3.5cm}{.25pt}\end{center}
\clearpage\endgroup
